% Generated draft: requires mathematical and source review.
\documentclass[11pt]{article}
\usepackage[margin=25mm]{geometry}
\usepackage{fontspec}
\setmainfont{Latin Modern Roman}
\usepackage{amsmath,amssymb,mathtools,unicode-math}
\setmathfont{Latin Modern Math}
\usepackage{enumitem,needspace,xurl,hyperref}
\setlength{\parindent}{0pt}
\setlength{\parskip}{5pt}
\setlength{\emergencystretch}{3em}
\newcommand{\N}{\mathbb{N}}
\newcommand{\Z}{\mathbb{Z}}
\newcommand{\Q}{\mathbb{Q}}
\newcommand{\R}{\mathbb{R}}
\newcommand{\C}{\mathbb{C}}
\newcommand{\F}{\mathbb{F}}
\newcommand{\A}{\mathbb{A}}
\newcommand{\PP}{\mathbb{P}}
\newcommand{\E}{\mathbb{E}}
\newcommand{\eps}{\varepsilon}
\newcommand{\dd}{\,\mathrm{d}}
\begin{document}
\section*{UnsolvedMath ID: 30001453}
\textbf{Source problem code:} catalogue entry
\subsection{Definitions and mathematical statement}
The supplied record requests a classification of doubly singular \((2,0)\)-irreducible partitions. The formulation below uses the usual interpretation of \((2,0)\): Hecke parameter \(q=-1\) over a field of characteristic zero, with irreducibility referring to a Specht module. These conventions and the intended base field require confirmation from the full source; they are not defined in the supplied fragment.

Fix a field \(K\) of characteristic zero. A partition of an integer \(n\geq 1\) is a finite sequence \(\lambda=(\lambda_1,\ldots,\lambda_r)\) with \(\lambda_1\geq\cdots\geq\lambda_r>0\) and \(\sum_{i=1}^{r}\lambda_i=n\). Its conjugate partition is defined by \(\lambda'_j=\#\{i:\lambda_i\geq j\}\). A partition is \(2\)-singular if some positive part occurs at least twice. It is \emph{doubly singular} if both \(\lambda\) and \(\lambda'\) are \(2\)-singular.

Let \(H_n(K,-1)\) be the type-A Iwahori--Hecke algebra with generators \(T_1,\ldots,T_{n-1}\), relations
\[
(T_i+1)^2=0,\qquad
T_iT_{i+1}T_i=T_{i+1}T_iT_{i+1},\qquad
T_iT_j=T_jT_i\quad (|i-j|>1),
\]
for all indices for which the displayed expressions are defined. Let \(S_K^\lambda\) denote its Specht module indexed by \(\lambda\). Under the interpretation above, \(\lambda\) is \((2,0)\)-irreducible if \(S_K^\lambda\) is a nonzero module with no nonzero proper submodule.

The requested classification, under these conventions and for the intended choice of \(K\), is to give an explicit necessary and sufficient combinatorial condition on \(\lambda\), valid for every \(n\geq1\), describing precisely the set
\[
\mathcal C_K=
\bigcup_{n\geq1}
\left\{\lambda\vdash n:
\lambda\text{ and }\lambda'\text{ are }2\text{-singular, and }
S_K^\lambda\text{ is irreducible}\right\}.
\]
The source verification describes this as the remaining characteristic-zero case; that description is supplied attribution, not an independently verified claim.

The original statement additionally reports a conjectural construction from rectangular partitions, meaning partitions \((a^b)\) with \(a,b\geq1\), by unspecified modifications. No modifications or exact criterion are supplied, so no equivalence with such a construction can be stated precisely here. The corrected statement deliberately omits that conjectural clause. Consequently, this draft records the classification target but does not reconstruct the missing conjecture.
\subsection{Short English statement}
Classify partitions whose rows and conjugate rows both contain repeated positive parts and whose associated type-A Hecke Specht modules are irreducible at parameter \(q=-1\) in characteristic zero. This uses an explicitly stated conventional interpretation of \((2,0)\); the intended field and the omitted conjectural modifications of rectangles require source confirmation. The supplied open status is dated curation, not an independently verified mathematical assessment.
\subsection{Sources}
\begin{itemize}
\item Local UnsolvedMath registry, permanent ID \(30001453\), title: \emph{Doubly Singular Irreducible Partitions in Type A}; exact source problem code: \texttt{catalogue entry}. The supplied record belongs to the Oberwolfach Reports Open Problems collection, with proposed year 2010; no individual proposer is identified.
\item Supplied report citation and literature source labeled ``original OWR report'': \emph{Combinatorial Representation Theory} (2010), Oberwolfach Reports, pp.\ 799--882, DOI \texttt{10.4171/owr/2010/15}. \url{https://doi.org/10.4171/owr/2010/15}. The supplied background describes Oberwolfach Reports as documenting mathematical workshops and their open-problem discussions. This report has not been independently consulted for this draft.
\item Supplied literature assessment, checked exactly \texttt{2026-08-21}: status recorded as open. It describes the extracted record as a truncated continuation with no complete problem statement and says that the unresolved remainder cannot be verified from the fragment. This is supplied dated curation, not an independently verified literature assessment.
\item The supplied statement status is \texttt{corrected\_verified}. Its accompanying verification says that the report identifies this classification as the remaining characteristic-zero case, and that the informal rectangle-modification clause was omitted because its precise conjectural criterion occurs only in a cited Fayers--Lyle paper. No title, publication details, or link for that paper are supplied, so its criterion is not reproduced here. This verification label and explanation are supplied registry metadata, not independent verification.
\end{itemize}
\end{document}
